\section*{Introducción}

El presente espacio curricular se fundamenta en la necesidad del desarrollo
sistemático de saberes relacionados con los principios técnicos y termodinámicos que dan fundamento teórico para la comprensión de las máquinas térmicas en general, compresores, motores de combustión interna, máquinas frigoríficas, turbinas de gas y vapor de forma contextualizada.
Termodinámica y Máquinas Térmicas posee formato asignatura. Profundiza
los saberes desarrollados en Física y Químicade 3° año. Articula también con Cálculo y Elementos de Máquina y Operaciones de Taller de Electromecánicas I en 5° año y finalmente los saberes generados en este espacio serán considerados en
Mantenimiento y Reparación de Equipos, en las Prácticas Profesionalizantes, y
principalmente en Instalaciones Industriales al realizar el diseño de instalaciones que involucren algún proceso termodinámico.

Es esencial que los estudiantes conozcan todo lo relacionado con las distintas
escalas de temperatura, prestando especial atención a la escala absoluta de orden térmico, introduciéndose así en el conocimiento de las transformaciones entre sistemas contemplando flujos circulantes, análisis de reversibilidad y el manejo de variables tales como volumen, presión y temperatura.
Es apropiado considerar el tratamiento de diversas técnicas de resolución de
problemáticas térmicas y transformaciones, a partir de consideraciones energéticas, siendo posible relacionar e interpretar unificadamente distintos fenómenos industriales.

Se generarán trabajos o proyectos interdisciplinarios que promuevan la
articulación de saberes con Operaciones de Taller Electromecánicas I, aprovechando el equipamiento tecnológico perteneciente a la institución educativa y con las Prácticas Profesionalizantes, que propician el aprendizaje práctico, mediante pasantías en las empresas del medio.

\section*{TEMAS}

\begin{enumerate}
    \item FORMAS DE ENERGÍA Y SUS TRANSFORMACIONES.
    \begin{itemize}
        \item Conocer los principios que
rigen las transformaciones
energéticas. 
        \begin{itemize}
            \item Interpretación de las distintas escalas de temperatura.
\item Reconocimiento de la fuerza impulsora (diferencia de
temperatura) que permite la transferencia de calor.
\item Diferenciación de los mecanismos de transferencia de
calor.(conducción, convección, radiación)
\item Análisis comparativo de dilataciones lineales, superficiales
y cúbicas. 
\item Identificación del Principio Cero de la Termodinámica, y su
aplicación en Calorimetría.
\item Elaboración de la relación existente entre balance térmico,
capacidad calorífica y calor específico.
\item Descripción de los fenómenos de dilatación y deformación
en materiales de uso industrial.
        \end{itemize}
    \end{itemize}
    
    
    \item PRINCIPIOS FUNDAMENTALES DE LA TERMODINÁMICA. 
    \begin{itemize}
        \item Analizar el trabajo en un
sistema cerrado adiabático.
\begin{itemize}
    \item Conceptualización de las características de trabajo
termodinámico.
\item Identificación de las condiciones de un sistema cerrado.
\item Caracterización del principio de la accesibilidad
adiabática.
\item Caracterización de los principios enunciados por Lord
Kelvin y Clausius.
\item Identificación del principio de conservación de la energía
entre estados de equilibrio.
\item Conocimiento de las interacciones másica, mecánica y
térmica de modo sistémico.
\item Representación gráfica de transformaciones isotérmicas,
isobáricas, isocóricas, adiabáticas, politrópicas.
\item Conceptualización de la energía interna como variable de
estado.
\item Análisis crítico del Primer Principio de la Termodinámica.
\item Caracterización de trabajo de expansión y circulación.
\item Representación gráfica de los distintos tipos de trabajos
termodinámicos.
\end{itemize}
        \item Conocer las causas y consecuencias del incremento de la Entropía en el Universo.
        \begin{itemize}
            \item  Caracterización de la Entropía de modo universal e industrial.
\item Explicitación de la irreversibilidad en procesos
termodinámicos.
\item Reconocimiento de teorías de cantidad de energía no
utilizable en los sistemas.
\item Análisis de los procesos de transformación de materia y
energía desde una perspectiva de creación-destrucción.
\item Análisis crítico del Segundo Principio de la termodinámica
desde la visión de Clausius y de lo desarrollado por
Kelvin-Planck.
\item Reconocimiento de la máquina ideal de Carnot
contemplando el rendimiento idealista y su reversibilidad.
\item Análisis de diagramas gráficos relacionados con procesos
entrópicos.
\item Formulación del Segundo Principio de la Termodinámica desde una perspectiva técnica - industrial. 
        \end{itemize} 
    \end{itemize}
    \item  MÁQUINAS DE COMBUSTIÓN INTERNA 
    \begin{itemize}
        \item Analizar críticamente la matriz tecnológica aplicada en las máquinas de combustión interna en el mundo energético actual.
    \begin{itemize}
        \item Conocimiento de los aportes sobre motores de
combustión interna de NikolausAugust Otto.
\item Reconocimiento de los tiempos o ciclos termodinámicos
en los motores de combustión interna.
\item Interpretación gráfica de las variables intervinientes en un
ciclo ideal de combustión interna.
\item Análisis comparativo entre los desarrollos tecnológicos de
Nikolaus Otto y Rudolf Karl Diesel, considerando
beneficios y cualidades de los ciclos de motores
generados por ellos.
\item Diferenciación de rendimientos de los distintos ciclos de
combustión interna.
\item Conceptualización de las fórmulas técnicas vinculadas
con las máquinas de combustión interna.
    \end{itemize}
    \end{itemize}
    \item CICLOS TERMODINÁMICOS DE VAPOR Y GAS.
    \begin{itemize}
        \item Reconocer y analizar las transformaciones termodinámicas en los ciclos de vapor.
        \begin{itemize}
            \item Caracterización de las transformaciones de un sistema
gaseoso.
\item Interpretación del fenómeno de vaporización, a través de
diagramas, representación, con usos de tablas de vapor.
\item Conceptualización de la Entalpia del líquido y del vapor
para el empleo tecnológico en máquinas refrigerantes.
\item Comprensión de los fenómenos de vaporización e
interpretación de las tablas de vapor.
\item Conocimiento de los principios de Carnot relacionados
con los ciclos de vapor.
\item Reconocimiento de las características de diagramas,
título de vapor, empleando los aportes técnicos de Mollier.
\item Caracterización de los principios desarrollados por William
John Rankine en los ciclos de vapor
\item Reconocimiento de los efectos de presión y temperatura
en los ciclos de vapor. 
        \end{itemize}
        \item Conocer los principios vinculados con los ciclos de gas. 
        \begin{itemize}
            \item Caracterización del funcionamiento de una turbomáquina
motora, con empleo de gas como fluido actuante.
\item Análisis comparativo de las propiedades técnico-térmicas
de los ciclos de vapor y gas.
\item Conceptualización de ciclos de gas de potencia y de
refrigeración según Brayton.
\item Exploración de la fluidez y la presión del aire en un ciclo
termodinámico.
\item Interpretación de las funciones tecnológicas específicas,
del compresor, turbina y otros dispositivos acoplados en la
generación eléctrica.
\item Diferenciación de eficiencias y rendimientos entre los
ciclos combinados. 
        \end{itemize}
    \end{itemize}
\end{enumerate}

\section*{Programa}

\subsection*{Primera parte:}  	Termodinámica. Calor y Temperatura. Definiciones. Unidades. Escalas termométricas fundamentales, Celsius, Fahrenheit y Kelvin. Conversión de unidades. Transmisión de calor por conducción, por radiación y convección. Formulas. Coeficientes de transmisión. 
Dilatación lineal, superficial y cubica, coeficientes. Calorimetría, definición, unidades, capacidad calorífica, calor especifico, unidades. Ecuación fundamental de la calorimetría. Trabajo termodinámico de expansión y circulación. Representación gráfica de los distintos tipos de trabajos. Definición principio cero de la termodinámica. Sistemas termodinámicos,  tipos, variables fundamentales diagrama P-V. Principio de conservación de la energía. Tipos de transformaciones termodinámicas, isotérmicas, isobáricas, isocóricas, adiabáticas, politrópicas. Graficas. Energía interna y su relación con la entropía. Procesos termodinámicos ideales e irreversibles. Segundo principio de la termodinámica. Relación con los enunciados de  Clausius y Kelvin-Planck. Maquina ideal de Carnot ciclo ideal. Diagramas de procesos entrópicos.
\subsection*{Segunda Parte:}  Maquinas térmicas. Definición, generalidades, usos. Motores de combustión interna. Descripción y funcionamiento.  Ciclos o tiempos termodinámicos. Ciclo Otto, rendimiento real e ideal. Fórmulas. Tiempos o ciclos termodinámicos en los motores de combustión interna. Ciclo Diesel. Rendimiento térmico. Interpretación gráfica de las variables intervinientes en un ciclo ideal de combustión interna. Ciclo semi Diesel. Turbinas de gas, ciclo Brayton. Descripción y funcionamiento.    Diferenciación de rendimientos de los distintos ciclos de combustión interna. Gráficas. Maquinas frigoríficas, clasificación, ciclo. 
Transformaciones de un sistema gaseoso. Fenómenos de vaporización. Tablas de vapor. Entalpia del líquido y del vapor en máquinas refrigerantes. Conocimiento de los Principios de Carnot para ciclos de vapor. Diagramas, título de vapor,  enunciado de Mollier.  Enunciado de William John Rankine. Reconocimiento de los efectos de presión y temperatura en los ciclos de vapor. Conocer los principios vinculados con los ciclos de gas. Funcionamiento de una turbomáquina motora a gas. Propiedades técnico-térmicas de los ciclos de vapor y gas. Ciclos de gas de potencia y de refrigeración según Brayton. Funciones del compresor, turbina y dispositivos acoplados en la generación eléctrica. Eficiencias y rendimientos entre los ciclos combinados.






\section*{Normas para la presentación de la carpeta de trabajos prácticos, teoría y acuerdos áulicos.}
\begin{enumerate}
    \item La portada debe ser la propuesta por el docente, en caso que se encuentra la portada de la escuela se utilizará la misma, el índice debe estar completo, con la numeración completa, en la misma debe estar el programa de la materia y la presentes pautas que deben estar firmadas por el alumno.
    \item Rotulos completos, realizados en regla, con todos los datos correspondientes y el formato propuesto por el docente, las hojas deben estar numeradas, letra técnica o imprenta prolija segun lo pactado po r el docente.
    \item La carpeta debe estar prolija, se tendrá en cuenta la ortografía, las hojas deben ser la misma sen cada trabajo práctico, al cambiar de trabajo pueden cambiar las hojas, se debe utilizar el mismo color de lapicera, los dibujos deben ser realizados en regla.
    \item El alumno deberá tener dos evaluaciones por informe, en total seis evaluaciones al año, si el alumno falta deberá traer el certificado correspondiente.
    \item \textbf{Trabajo en clase, trabajo en grupo, pizarrón, actitudinal:} El alumno deberá tener un vocabulario adecuado en el aula, respeto hacia el docente y hacia sus compañeros y las normas de convivencia.
    \item Carpeta completa y entrega en termino.
    \item Asistencia (80\%).
\end{enumerate}
\vspace{1cm}
\begin{center}
Firma Alumno $\dots\dots\dots\dots\dots\dots\dots\dots\dots\dots$

\end{center}

\tableofcontents


